% =============================================================================
% MOLAMBUDOS —— 患者1,260的日记 (CHINESE EDITION 简体中文)
% Template Victoria-Regia by @jancarauma — XeLaTeX + xeCJK variant
% 84 fragments · 216 reading paths · 5 parts (SPEC-935-R356; contagens atualizadas no R401)
% Mirror of main.tex / en/main_en.tex
% Compile: xelatex -interaction=nonstopmode -halt-on-error zh/main_zh.tex
% =============================================================================

\documentclass[14pt,openany]{extbook}
\newif\ifmolambudoskdpprint
\ifdefined\KDPPRINT
  \molambudoskdpprinttrue
\else
  \molambudoskdpprintfalse
\fi
\newcommand{\goldenratio}{1.618}

% --- Book metadata ---
\newcommand{\authorname}{Marcelo Dias de Carvalho Filho}
\newcommand{\coauthorname}{Marcelo Claro}
\newcommand{\translatorname}{}
\newcommand{\booktitle}{Molambudos —— 患者1,260的日记}
\newcommand{\subtitle}{碎片式叙事}
\newcommand{\publisher}{OpenCode Books}
\newcommand{\publishercity}{Crateús}
\newcommand{\editionyear}{2026}
\newcommand{\isbn}{9798189170492}

% --- Template packages (XeLaTeX variant) + customizations ---
\usepackage{misc/options_zh}
\usepackage{float}

% --- Figure directory (sepia photographs of the diary) ---
\graphicspath{{figures/}}

% ============================================================================
% Molambudos-specific customizations
% ============================================================================

% Conditional hypertext — recognizes the reading path
% \ifpath{origin1}{origin2}{text if came from 1}{text if came from 2}
\newcommand{\ifpath}[4]{%
  \textit{如果您来自#1:} #3\par
  \vspace{0.3em}
  \textit{如果您来自#2:} #4\par
  \vspace{0.3em}
  \textit{（两条路径通向同一处。实体无需知道您来自哪里。它只需要您已经抵达。）}%
}

% Phantom route — links to a fragment that does not exist.
\newcommand{\rotaphantom}[1]{%
  \textbf{#1}{\footnotesize(p.~---)}%
}

% Molambudos color palette
\definecolor{sangue}{HTML}{8B0000}
\definecolor{olho}{HTML}{FFD700}
\definecolor{mem}{HTML}{8B4513}
\definecolor{doc}{HTML}{2F4F4F}
\definecolor{cont}{HTML}{8B0000}
\definecolor{luc}{HTML}{4B0082}
% Immersive sepia palette (page background + ink) — sensory appeal.
\definecolor{sepiabg}{HTML}{F2E8CF}
\definecolor{sepiaink}{HTML}{2B2117}

% Hyperlink setup (internal hyperlinks between fragments)
\ifmolambudoskdpprint
  \hypersetup{
    colorlinks=false,
    hidelinks
  }
\else
  \hypersetup{
    colorlinks=true,
    linkcolor=linkred,
    citecolor=linkred,
    urlcolor=linkred
  }
\fi

% Part opener — title + thematic roundel (icone/N.png)
% on a single clean page. #1 displayed title, #2 short name (TOC), #3 icon.
\newcommand{\partopener}[3]{%
  \cleardoublepage
  \thispagestyle{empty}%
  \phantomsection\addcontentsline{toc}{part}{#2}%
  \markboth{}{}%
  \begin{tikzpicture}[remember picture,overlay]
    \fill[black] (current page.south west) rectangle (current page.north east);
    \node[anchor=center] at (current page.center){%
      \begin{tabular}{c}
        {\bfseries\LARGE\color{customgold} #1}\\[2.4em]
        \includegraphics[width=0.58\textwidth]{icone/#3}\\
      \end{tabular}};
  \end{tikzpicture}%
  \null\vfill\clearpage}

% Reading route — link to a fragment + its page number
\ifmolambudoskdpprint
  \newcommand{\rota}[1]{\textbf{#1}\,{\footnotesize(p.~\pageref{frag:#1})}}
\else
  \newcommand{\rota}[1]{\hyperlink{#1}{\textbf{#1}}\,{\footnotesize(p.~\pageref{frag:#1})}}
\fi
\newcommand{\rotasep}{\quad\textbullet\quad}

% Marginalia of the Contamination fragments
\setlength{\marginparwidth}{0.75in}
\setlength{\marginparsep}{6pt}
\ifmolambudoskdpprint
  % KDP print mode: margin notes suppressed
  \newcommand{\margnota}[1]{}
\else
  \newcommand{\margnota}[1]{\marginpar{\raggedright\footnotesize\itshape\color{sangue}#1}}
\fi

% Editor's Note — paratextual annotation [编者按]
\newcommand{\NE}[1]{\footnote{\footnotesize\textbf{[编者按]} #1}}

% Fragment counter
\newcounter{fragctr}
\newcommand{\fragdef}[1]{%
  \refstepcounter{fragctr}%
  \hypertarget{#1}{}%
  \label{frag:#1}%
}

\begin{document}

% 封面和全部前置页共用一套罗马页码；正文阿拉伯页码开始前关闭自动页锚。
\hypersetup{pageanchor=false}
\frontmatter

\input{zh/frontmatter/titlepage}

\input{zh/frontmatter/aviso_ao_leitor}

\ifdefined\MolambudosKDPNoFolios
  \makeatletter
  \let\ps@plain\ps@empty
  \makeatother
  \pagestyle{empty}
\fi
\ifdefined\MolambudosKDPSafeFolios
  % KDP print mode with preserved physical navigation
  \setlength{\footskip}{6mm}
  \fancypagestyle{molambudos-kdp-folios}{%
    \fancyhf{}%
    \renewcommand{\headrulewidth}{0pt}%
    \renewcommand{\footrulewidth}{0pt}%
    \fancyfoot[C]{\normalfont\footnotesize\thepage}%
  }%
  \fancypagestyle{plain}{%
    \fancyhf{}%
    \renewcommand{\headrulewidth}{0pt}%
    \renewcommand{\footrulewidth}{0pt}%
    \fancyfoot[C]{\normalfont\footnotesize\thepage}%
  }%
  \pagestyle{molambudos-kdp-folios}
\fi

% Sepia background (parchment) + dark sepia ink
\pagecolor{sepiabg}
\color{sepiaink}

\input{zh/frontmatter/epigrafe}
\input{zh/frontmatter/cuidado}
\input{zh/frontmatter/copyrightpage}
\input{zh/frontmatter/ficha_catalografica}
\cleardoublepage
\input{zh/frontmatter/nota_arquivista}
\input{zh/frontmatter/como_ler}
\input{zh/frontmatter/tocpage}

% ===== 前置导航 =====
\pagestyle{fancy}
\ifdefined\MolambudosKDPNoFolios
  \pagestyle{empty}
\fi
\ifdefined\MolambudosKDPSafeFolios
  \pagestyle{molambudos-kdp-folios}
\fi

% ============================================
% NAVIGATION SECTION — Non-Linear Reading
% ============================================
\chapter*{导航——阅读路线}
\addcontentsline{toc}{chapter}{导航——阅读路线}

\noindent\textit{本档案包含按5个部分组织的78个碎片。}
\textit{每个碎片都是独立的，可以按任意顺序阅读。}
\textit{每个碎片末尾的箭头 ↪ 表示可能的阅读路径。}
\textit{推荐的线性阅读遵循以下各部分顺序。}

\vspace{1em}

\section*{碎片索引}

% Compact fraglink for the index
\ifmolambudoskdpprint
  \newcommand{\fraglink}[3]{%
    \par\noindent%
    \textbf{#1} --- #2\,{\footnotesize(p.~\pageref{frag:#1})}%
    \ifx&#3&\else\hfill {\itshape\textcolor{textgray}{#3}}\fi%
    \ifdefined\MolambudosCompactIndex\par\vspace{0pt}\else\par\vspace{0.1em}\fi}
\else
  \newcommand{\fraglink}[3]{%
    \par\noindent%
    \hyperlink{#1}{\textbf{#1}} --- #2\,{\footnotesize(p.~\pageref{frag:#1})}%
    \ifx&#3&\else\hfill {\itshape\textcolor{textgray}{#3}}\fi%
    \par\vspace{0.1em}}
\fi

\subsection*{第一部 —— 塞尔唐（9个碎片）}
{\setlength{\parskip}{0.15em}
\fraglink{MEM-01}{终结前的塞尔唐}{1914}
\fraglink{MEM-02}{逃荒者的跋涉}{1915}
\fraglink{MEM-03}{父亲的死}{1915}
\fraglink{MEM-04}{乱葬坑}{1915}
\fraglink{MEM-05}{黄眼}{1915}
\fraglink{MEM-06}{政府牲畜圈}{1915}
\fraglink{MEM-07}{被带走的母亲}{1915}
\fraglink{DOC-17}{一个士兵的信}{1915}
\fraglink{DOC-18}{圈栏官方报告}{1916}
}

\subsection*{第二部 —— 收容院（16个碎片）}
{\setlength{\parskip}{0.15em}
\fraglink{MEM-08}{抵达收容院}{1917}
\fraglink{DOC-02}{患者1,260病历}{1917}
\fraglink{MEM-09}{最初几年}{1920--1930}
\fraglink{MEM-10}{第一起谋杀}{1922}
\fraglink{MEM-11}{隔离}{1923--1945}
\fraglink{MEM-17}{电休克}{1937}
\fraglink{MEM-18}{逃跑的尝试}{1943}
\fraglink{MEM-12}{裹尸布撕裂者}{1945}
\fraglink{MEM-21}{报纸}{1948}
\fraglink{MEM-19}{访客}{1950}
\fraglink{DOC-10}{ECT记录}{1937}
\fraglink{MEM-22}{声音}{1952}
\fraglink{MEM-23}{灰雪}{1960}
\fraglink{MEM-20}{牢房邻居之死}{1963}
\fraglink{MEM-24}{哭泣的护士}{1964}
\fraglink{MEM-25}{镜中之死}{1970}
}

\subsection*{第三部 —— 奥利维拉日记\\\hspace*{2.5em}与鉴定报告（25个碎片）}
{\setlength{\parskip}{0.15em}
\fraglink{DOC-01}{掘尸}{1979}
\fraglink{DOC-03}{尸检报告}{1979}
\fraglink{DOC-04}{会谈报告I}{1979}
\fraglink{DOC-05}{会谈报告II}{1979}
\fraglink{DOC-06}{会谈报告III}{1979}
\fraglink{DOC-07}{精神科诊断（1979）}{1979}
\fraglink{MEM-13}{新医生}{1979}
\fraglink{MEM-14}{日记浮现}{1979}
\fraglink{MEM-15}{实体开口}{1979}
\fraglink{MEM-16}{收容院之死}{1979}
\fraglink{MEM-26}{最后一页}{1979}
\fraglink{DOC-11}{奥利维拉日记I}{1979}
\fraglink{DOC-12}{奥利维拉日记II}{1979}
\fraglink{DOC-13}{奥利维拉日记III}{1979}
\fraglink{DOC-14}{奥利维拉日记IV}{1979}
\fraglink{DOC-15}{奥利维拉日记V}{1979}
\fraglink{DOC-16}{奥利维拉日记VI}{1979}
\fraglink{DOC-19}{巴西大屠杀}{达妮埃拉·阿尔贝克斯}
\fraglink{DOC-20}{保管链文书}{2026}
\fraglink{DOC-21}{埃托尔·奥利维拉的信（未寄出）}{1979}
\fraglink{DOC-22}{内部备忘录第47/2026号}{2026}
\fraglink{DOC-23}{档案馆事件记录}{2026}
\fraglink{DOC-24}{保护技术意见书}{2026}
\fraglink{DOC-25}{黑布男人}{2026}
\fraglink{DOC-26}{奥利维拉的笔记——机制}{1979}
}

\subsection*{第四部 —— 卢西亚调查（20个碎片）}
{\setlength{\parskip}{0.15em}
\fraglink{DOC-08}{精神科再诊断（2026）}{卢西亚·门德斯}
\fraglink{DOC-09}{患者1,263评估量表}{临床工具}
\fraglink{DOC-27}{起源（1853）}{若泽·伊纳西奥神父}
\fraglink{LUC-01}{邀请}{2026}
\fraglink{LUC-02}{档案}{2026}
\fraglink{LUC-03}{废墟}{2026}
\fraglink{LUC-04}{照片}{2026}
\fraglink{LUC-05}{卢西亚的梦}{2026}
\fraglink{LUC-Contraprova}{反证}{2026}
\fraglink{LUC-Entrevista}{访谈}{2026}
\fraglink{LUC-Escolha}{抉择}{2026}
\fraglink{LUC-06}{失踪}{2026}
\fraglink{LUC-07}{奥利维拉的踪迹}{2026}
\fraglink{LUC-08}{卢西亚最后一次会谈}{2026}
\fraglink{LUC-09}{奥利维拉的掘尸}{2026}
\fraglink{LUC-10}{卢西亚的鉴定书}{2026}
\fraglink{LUC-11}{数字日记}{2026}
\fraglink{LUC-12}{最后一张照片}{2026}
\fraglink{LUC-13}{诊断之夜}{2026}
\fraglink{LUC-14}{临床自评（DSM-5-TR）}{2026}
}

\subsection*{第五部 —— 污染（14个碎片＋尾声）}
{\setlength{\parskip}{0.15em}
\fraglink{CONT-01}{您打开了日记}{}
\fraglink{CONT-02}{您当时就在那里}{}
\fraglink{CONT-03}{您闻到了气味}{}
\fraglink{CONT-04}{您开始做梦}{}
\fraglink{CONT-05}{临床评估报告}{}
\fraglink{CONT-06}{您已经读过这一行三遍}{}
\fraglink{CONT-07}{最终诊断}{}
\fraglink{CONT-08}{您以前读过这一切}{}
\fraglink{CONT-09}{镜子}{}
\fraglink{CONT-10}{印记}{}
\fraglink{CONT-11}{数字}{}
\fraglink{CONT-12}{3:14}{}
\fraglink{CONT-13}{队列}{}
\fraglink{MEM-27}{下一页}{若阿金，1907--1979}
}

\newpage

% ============================================
% NARRATIVE GRAPHS — Visual Navigation of the Work
% ============================================
\chapter*{视觉导航——档案地图}
\addcontentsline{toc}{chapter}{视觉导航——地图}

\noindent\textit{在接下来的几页中，三张地图帮助您为阅读定向。}
\textit{每张各有用途：探索、理解，或跟随。}
\textit{没有一张是必须的。所有都是可能的路径。}

\vspace{0.6em}
\noindent\textbf{地图1} —— \textbf{路径图}展示所有碎片及其连接。用于\textit{非线性导航}：选择一个碎片，跟随箭头。

\noindent\textbf{地图2} —— \textbf{情节图}揭示作品在三个幕中的隐秘结构。其中蕴含着原因、传播，以及连接四位宿主——包括您——的循环。

\noindent\textbf{地图3} —— \textbf{线性图}是传统路线，从塞尔唐（1914）到尾声。推荐给喜欢从头开始的人。

\vspace{0.8em}
\noindent{\small\textit{路径的选择不会改变目的地：所有路线都汇聚于同一个循环。}}

%% --- GRAPH 1: PATH MAP (rotated 90°) ---
\newpage
\thispagestyle{empty}
\begingroup
\vbox to \textheight{%
  \vfill
  \centering
  \setlength{\fboxsep}{3pt}%
  \setlength{\fboxrule}{0.6pt}%
  \fbox{\includegraphics[angle=90,width=0.98\textwidth,height=0.90\textheight,keepaspectratio]{misc/grafo_narrativo.png}}%
  \vfill
}%
\endgroup
\newpage

\section*{地图1 —— 阅读图：碎片之间的路径}

\noindent 本图表示\textbf{78个碎片}（每个碎片是一个节点）及其\textbf{180条连接}（连接它们的边）。每个碎片在文本末尾通过链接 \texttt{↪} 指向其他碎片，并附有页码。

\noindent 一个碎片的连接越多，它在叙事网络中的位置就越核心。连接最多的碎片是：

\begin{itemize}
\itemsep0em
\item \textbf{MEM-04}（乱葬坑）—— 引发一切的相遇
\item \textbf{CONT-04}（您开始做梦）—— 污染抵达读者
\item \textbf{DOC-01}（掘尸）—— 日记从过去浮现
\end{itemize}

\vspace{0.5em}
\noindent\textbf{图例:}
\textcolor{mem}{MEM（棕色）} —— 患者1,260的记忆 \quad
\textcolor{doc}{DOC（灰色）} —— 文件与报告 \quad
\textcolor{cont}{CONT（红色）} —— 读者的污染 \quad
\textcolor{luc}{LUC（靛蓝）} —— 法证调查

%% --- GRAPH 2: PLOT STRUCTURE ---
\newpage
\thispagestyle{empty}
\begingroup
\vbox to \textheight{%
  \vfill
  \centering
  \setlength{\fboxsep}{3pt}%
  \setlength{\fboxrule}{0.6pt}%
  \fbox{\includegraphics[angle=90,width=0.98\textwidth,height=0.90\textheight,keepaspectratio]{misc/grafo_enredo.png}}%
  \vfill
}%
\endgroup
\newpage

\section*{地图2 —— 情节图：三幕架构}

\begingroup\raggedright

\noindent 与前一张地图（展示\textit{连接}）不同，此图揭示作品的\textit{因果结构}——支撑78个碎片的叙事骨架。

\vspace{1em}
\noindent\textbf{三幕:}

\vspace{0.4em}
\noindent\textcolor{mem}{\textbf{第一幕 —— 创伤}} \textit{（棕色）}:
1853年诞生实体的大旱，1915年将若阿金带往乱葬坑的大旱，
一个八岁男孩与永无止境的饥饿订下的契约。\\[0.35em]
\noindent\textcolor{doc}{\textbf{第二幕 —— 制度化}} \textit{（灰色）}:
收容院中的62年，作证的日记，将其归还世界的掘尸，
以及那个因为阅读而成为下一环的医生。\\[0.35em]
\noindent\textcolor{cont}{\textbf{第三幕 —— 循环}} \textit{（红色）}:
卢西亚的调查，发现恐怖是一个循环，
读者的污染——以及最终的领悟：
\textit{循环随每位读者重新开始}。

\vspace{0.9em}
\noindent\textbf{宿主队列:} 虚线边
(\textcolor{sangue}{-\,-\,-}) 连接那些一个接一个携带怪物的人
—— 若阿金（1,260）、奥利维拉（1,261）、卢西亚（1,262）和您（1,263）——，
构成一条始于1853年、没有宣告终点的链条。

\vspace{0.5em}
\noindent\textit{注: 碎片DOC-27（起源，1853）和LUC-13（诊断之夜）是在卢西亚研究期间加入档案的。
DOC-27是全档案最古老的文献（1853年），揭示了实体的起源；LUC-13记录了卢西亚从调查者
变成宿主的那个确切时刻。}

\endgroup

%% --- GRAPH 3: LINEAR ROUTE (rotated 90°) ---
\newpage
\thispagestyle{empty}
\begingroup
\vbox to \textheight{%
  \vfill
  \centering
  \setlength{\fboxsep}{3pt}%
  \setlength{\fboxrule}{0.6pt}%
  \fbox{\includegraphics[angle=90,width=0.98\textwidth,height=0.90\textheight,keepaspectratio]{misc/grafo_leitura_linear.png}}%
  \vfill
}%
\endgroup
\newpage

\section*{地图3 —— 线性阅读图：传统路线}

\noindent 对偏好顺序阅读的读者，此图按推荐的编年叙事顺序组织\textbf{78个碎片}：
从塞尔唐（1914）到尾声——穿越收容院、奥利维拉的日记和卢西亚的调查，直到最终的污染。

\vspace{0.8em}

\noindent\textbf{线性路线:}
\textcolor{mem}{塞尔唐（9）} $\to$ \textcolor{doc}{收容院（16）} $\to$
\textcolor{doc}{奥利维拉日记（18）} $\to$
\textcolor{luc}{卢西亚调查（21）} $\to$
\textcolor{cont}{污染（13）} $\to$
\textcolor{olho}{MEM-27（1）}

\vspace{0.3em}
\noindent{\small $\bigstar$ = 标星的碎片是基本最短路线上必须阅读的。}

\vspace{0.7em}
\noindent\textit{阅读提示:} 如果您是第一次进入档案，
从塞尔唐开始。若阿金的记忆为之后的一切铺好了地基。如果您偏爱谜团，从卢西亚调查
（第四部）开始，像发掘一段过去那样发现更早的碎片。

\vspace{0.3em}
\noindent{\small \textcolor{sangue}{\textbf{警告:}} 在任何路线上，第五部（污染）都必须最后阅读。}
\textit{它不是关于故事的。它是关于您的。}

\newpage

% ============================================
% LINEAR READING — 78 FRAGMENTS IN 5 PARTS
% ============================================

%% ============================================
%% PART 1 — THE SERTÃO (1915--1928)
%% ============================================
\partopener{塞尔唐（1914--1917）}{塞尔唐}{2.png}

% 第一个碎片从阿拉伯数字第1页开始；此前不得插入 \newpage。
\mainmatter
\hypersetup{pageanchor=true}
\fragdef{MEM-01}\input{zh/fragmentos/mem/MEM-01}
\newpage\fragdef{MEM-02}\input{zh/fragmentos/mem/MEM-02}
\newpage\fragdef{MEM-03}\input{zh/fragmentos/mem/MEM-03}
\newpage\fragdef{MEM-04}\input{zh/fragmentos/mem/MEM-04}
\newpage\fragdef{MEM-05}\input{zh/fragmentos/mem/MEM-05}
\newpage\fragdef{MEM-06}\input{zh/fragmentos/mem/MEM-06}
\newpage\fragdef{MEM-07}\input{zh/fragmentos/mem/MEM-07}
\newpage\fragdef{DOC-17}\input{zh/fragmentos/doc/DOC-17}
\newpage\fragdef{DOC-18}\input{zh/fragmentos/doc/DOC-18}

%% ============================================
%% PART 2 — THE COLONY (1928--1979)
%% ============================================
% TODO-R356: translate zh/fragmentos/{mem,doc} PART 2, then uncomment
\partopener{收容院（1917--1979）}{收容院}{4.png}

\newpage\fragdef{MEM-08}\input{zh/fragmentos/mem/MEM-08}
\newpage\fragdef{DOC-02}\input{zh/fragmentos/doc/DOC-02}
\newpage\fragdef{MEM-09}\input{zh/fragmentos/mem/MEM-09}
\newpage\fragdef{MEM-10}\input{zh/fragmentos/mem/MEM-10}
\newpage\fragdef{MEM-11}\input{zh/fragmentos/mem/MEM-11}
\newpage\fragdef{MEM-17}\input{zh/fragmentos/mem/MEM-17}
\newpage\fragdef{MEM-18}\input{zh/fragmentos/mem/MEM-18}
\newpage\fragdef{MEM-12}\input{zh/fragmentos/mem/MEM-12}
\newpage\fragdef{MEM-21}\input{zh/fragmentos/mem/MEM-21}
\newpage\fragdef{MEM-19}\input{zh/fragmentos/mem/MEM-19}
\newpage\fragdef{DOC-10}\input{zh/fragmentos/doc/DOC-10}
\newpage\fragdef{MEM-22}\input{zh/fragmentos/mem/MEM-22}
\newpage\fragdef{MEM-23}\input{zh/fragmentos/mem/MEM-23}
\newpage\fragdef{MEM-20}\input{zh/fragmentos/mem/MEM-20}
\newpage\fragdef{MEM-24}\input{zh/fragmentos/mem/MEM-24}
\newpage\fragdef{MEM-25}\input{zh/fragmentos/mem/MEM-25}

%% ============================================
%% PART 3 — OLIVEIRA'S DIARY + REPORTS (1979)
%% ============================================
% TODO-R356: translate zh/fragmentos PART 3, then uncomment
\ifmolambudoskdpprint
  \partopener{\shortstack{奥利维拉日记\\与鉴定报告（1979）}}{奥利维拉日记与鉴定报告}{1.png}
\else
  \partopener{奥利维拉日记与鉴定报告（1979）}{奥利维拉日记与鉴定报告}{1.png}
\fi

\vspace{0.5em}
\noindent{\footnotesize\textit{档案员按语：从此刻起，主要的声音不再是若阿金的，而是那个研究他的人。以下的文件由医生、检验人员和档案员写成。语言变了。目光变了。但日记还是同一本。}}

\newpage\fragdef{DOC-01}\input{zh/fragmentos/doc/DOC-01}
\newpage\fragdef{DOC-03}\input{zh/fragmentos/doc/DOC-03}
\newpage\fragdef{DOC-04}\input{zh/fragmentos/doc/DOC-04}
\newpage\fragdef{DOC-05}\input{zh/fragmentos/doc/DOC-05}
\newpage\fragdef{DOC-06}\input{zh/fragmentos/doc/DOC-06}
\newpage\fragdef{DOC-07}\input{zh/fragmentos/doc/DOC-07}
\newpage\fragdef{MEM-13}\input{zh/fragmentos/mem/MEM-13}
\newpage\fragdef{MEM-14}\input{zh/fragmentos/mem/MEM-14}
\newpage\fragdef{MEM-15}\input{zh/fragmentos/mem/MEM-15}
\newpage\fragdef{MEM-16}\input{zh/fragmentos/mem/MEM-16}
\newpage\fragdef{MEM-26}\input{zh/fragmentos/mem/MEM-26}
\newpage\fragdef{DOC-11}\input{zh/fragmentos/doc/DOC-11}
\newpage\fragdef{DOC-12}\input{zh/fragmentos/doc/DOC-12}
\newpage\fragdef{DOC-13}\input{zh/fragmentos/doc/DOC-13}
\newpage\fragdef{DOC-14}\input{zh/fragmentos/doc/DOC-14}
\newpage\fragdef{DOC-15}\input{zh/fragmentos/doc/DOC-15}
\newpage\fragdef{DOC-16}\input{zh/fragmentos/doc/DOC-16}
\newpage\fragdef{DOC-19}\input{zh/fragmentos/doc/DOC-19}
\newpage\fragdef{DOC-20}\input{zh/fragmentos/doc/DOC-20}
\newpage\fragdef{DOC-21}\input{zh/fragmentos/doc/DOC-21}
\newpage\fragdef{DOC-22}\input{zh/fragmentos/doc/DOC-22}
\newpage\fragdef{DOC-23}\input{zh/fragmentos/doc/DOC-23}
\newpage\fragdef{DOC-24}\input{zh/fragmentos/doc/DOC-24}
\newpage\fragdef{DOC-25}\input{zh/fragmentos/doc/DOC-25}
\newpage\fragdef{DOC-26}\input{zh/fragmentos/doc/DOC-26}

%% ============================================
%% PART 4 — LÚCIA'S INVESTIGATION (2026)
%% ============================================
% TODO-R356: translate zh/fragmentos PART 4, then uncomment
\partopener{卢西亚调查（2026）}{卢西亚调查}{3.png}

\vspace{0.5em}
\noindent{\footnotesize\textit{档案员按语：本部分的碎片由卢西亚·门德斯医生在2026年研究期间整理。她的笔记以斜体写在原始文件的页边。档案中最古老的文献——来自1853年——由她发现并附入案卷。}}

\newpage\fragdef{DOC-08}\input{zh/fragmentos/doc/DOC-08}
\newpage\fragdef{DOC-09}\input{zh/fragmentos/doc/DOC-09}
\newpage\fragdef{DOC-27}\input{zh/fragmentos/doc/DOC-27}
\newpage\fragdef{LUC-01}\input{zh/fragmentos/luc/LUC-01}
\newpage\fragdef{LUC-02}\input{zh/fragmentos/luc/LUC-02}
\newpage\fragdef{LUC-03}\input{zh/fragmentos/luc/LUC-03}
\newpage\fragdef{LUC-04}\input{zh/fragmentos/luc/LUC-04}
\newpage\fragdef{LUC-05}\input{zh/fragmentos/luc/LUC-05}
\newpage\fragdef{LUC-Contraprova}\input{zh/fragmentos/luc/LUC-Contraprova}
\newpage\fragdef{LUC-Entrevista}\input{zh/fragmentos/luc/LUC-Entrevista}
\newpage\fragdef{LUC-Escolha}\input{zh/fragmentos/luc/LUC-Escolha}
\newpage\fragdef{LUC-06}\input{zh/fragmentos/luc/LUC-06}
\newpage\fragdef{LUC-07}\input{zh/fragmentos/luc/LUC-07}
\newpage\fragdef{LUC-08}\input{zh/fragmentos/luc/LUC-08}
\newpage\fragdef{LUC-09}\input{zh/fragmentos/luc/LUC-09}
\newpage\fragdef{LUC-10}\input{zh/fragmentos/luc/LUC-10}
\newpage\fragdef{LUC-11}\input{zh/fragmentos/luc/LUC-11}
\newpage\fragdef{LUC-12}\input{zh/fragmentos/luc/LUC-12}
\newpage\fragdef{LUC-13}\input{zh/fragmentos/luc/LUC-13}
\newpage\fragdef{LUC-14}\input{zh/fragmentos/luc/LUC-14}

%% ============================================
%% PART 5 — CONTAMINATION (YOU)
%% ============================================
% TODO-R356: translate zh/fragmentos/cont PART 5, then uncomment
\partopener{污染（您）}{污染}{5.png}

\newpage\fragdef{CONT-01}\input{zh/fragmentos/cont/CONT-01}
\newpage\fragdef{CONT-02}\input{zh/fragmentos/cont/CONT-02}
\newpage\fragdef{CONT-03}\input{zh/fragmentos/cont/CONT-03}
\newpage\fragdef{CONT-04}\input{zh/fragmentos/cont/CONT-04}
\newpage\fragdef{CONT-05}\input{zh/fragmentos/cont/CONT-05}
\newpage\fragdef{CONT-06}\input{zh/fragmentos/cont/CONT-06}
\newpage\fragdef{CONT-07}\input{zh/fragmentos/cont/CONT-07}
\newpage\fragdef{CONT-08}\input{zh/fragmentos/cont/CONT-08}
\newpage\fragdef{CONT-09}\input{zh/fragmentos/cont/CONT-09}
\newpage\fragdef{CONT-10}\input{zh/fragmentos/cont/CONT-10}
\newpage\fragdef{CONT-11}\input{zh/fragmentos/cont/CONT-11}
\newpage\fragdef{CONT-12}\input{zh/fragmentos/cont/CONT-12}
\newpage\fragdef{CONT-13}\input{zh/fragmentos/cont/CONT-13}

\newpage

%% ============================================
%% MEM-27 — JOAQUIM'S EPILOGUE
%% ============================================
% TODO-R356: translate zh/fragmentos/mem/MEM-27, then uncomment
\newpage\fragdef{MEM-27}\input{zh/fragmentos/mem/MEM-27}

\newpage

% ============================================
% EPILOGUE
% ============================================
\fragdef{EPI-01}
\chapter*{尾声}
\addcontentsline{toc}{chapter}{尾声}

\begin{center}
\vspace{2em}
\textit{循环随每位读者重新开始。}

\vspace{2em}
\rule{0.3\textwidth}{0.5pt}

\vspace{2em}
\textit{档案在此结束。}
\textit{但循环不会。}

\vspace{1.5em}
\rule{0.3\textwidth}{0.5pt}

\vspace{1em}
\textit{欢迎来到档案。}

\vspace{0.5em}
\textit{气味已经开始了吗？}

\vspace{1.5em}
--- 卢西亚·门德斯，患者1,262

\vspace{0.5em}
\rule{0.3\textwidth}{0.5pt}

\vspace{0.5em}
\noindent\textit{档案注：米纳斯吉拉斯州公共档案馆馆长罗纳尔多·托雷斯医生于2026年10月申请休假。医疗表格记录“持续性失眠、前额头痛和耳鸣”。没有提到眼睛。没有提到气味。协议日志中的最后一条，以他本人的笔迹写道：“交还患者1,263，待其出现。”其下，以检验人员无法辨认的笔迹写着：“等待。”}

\vspace{0.3em}
\noindent\textit{夜里，嗡鸣声像猫头鹰的歌一样起落——细弱、高亢，像一个女人在空档案馆的走廊里歌唱。没有人打开门去查看。}
\end{center}

\newpage
\thispagestyle{empty}
\null
\vfill
\begin{center}
\vspace*{\fill}

\textit{这本书包含78个碎片。}
\textit{如果您全部读完了，循环就在您体内完成了。}

\vspace{2em}
\rule{0.2\textwidth}{0.5pt}

\vspace{2em}
\textit{它现在就在这里。}
\textit{在空白的书页上。}
\textit{等待下一位读者。}

\vspace{2em}
\textcolor{olho}{\large \textbf{1260}}
\end{center}

\vfill
\clearpage

% ============================================
% EPILOGUE — SECOND PART (POST-READING)
% ============================================
\thispagestyle{empty}
\null
\vfill
\begin{center}

\vspace*{\fill}

\noindent\textit{\textcolor{sangue}{\textbf{读后指示}}}

\vspace{1.5em}

\noindent\textit{现在合上这本书。}

\noindent\textit{先别放回书架。把它留在桌上。}

\vspace{0.8em}

\noindent\textit{数62秒。}

\noindent\textit{什么也不想。只是数。}

\vspace{0.8em}

\noindent\textit{现在重新翻开这本书。翻到第62页。}

\noindent\textit{读第一行。}

\vspace{1em}

\noindent\textcolor{sangue}{\rule{0.15\textwidth}{0.5pt}}

\vspace{1em}

\noindent\textit{如果那一行写着“六十二”: 循环已经在您身上闭合。那个标记了若阿金一生的数字——62年、62天、62秒——现在属于您了。记住这个数字。它会在您意想不到的地方出现。}

\noindent\textit{如果那一行写着别的东西: 您以错误的顺序读了这本书。或者正确的顺序。无从知晓。第62页会随着打开它的人而改变。}

\vspace{1.5em}

\noindent\textcolor{sangue}{\rule{0.15\textwidth}{0.5pt}}

\vspace{1.5em}

\noindent\textit{在最终收好这本书之前:}

\vspace{0.5em}

\noindent\textit{--- 用右手抚过第126页。感受纸的纹理。}
\noindent\textit{--- 把手凑到鼻前。吸气。}
\noindent\textit{--- 如果甜味比第一次更浓，诊断已被确认。}
\noindent\textit{--- 如果气味更淡，您还处于潜伏期。}

\vspace{1.5em}

\noindent\textcolor{sangue}{\rule{0.15\textwidth}{0.5pt}}

\vspace{1.5em}

\noindent\textit{患者1,263的档案已经填好。}

\noindent\textit{“姓名”一栏是空白的。}

\noindent\textit{“污染日期”一栏填着您读到这一行的日期。}

\noindent\textit{欢迎来到档案。}

\noindent\textit{队列并没有变短。}

\vspace{1.5em}

\noindent\textit{--- 档案员}
\vspace{0.3em}
\noindent\textit{--- 患者1,262}
\vspace{0.3em}
\noindent\textit{--- 若阿金留下的残片}

\vspace{2em}
\textcolor{olho}{\large \textbf{1,263}}

\end{center}
\vfill

% ===== Editorial Appendix =====
\input{zh/frontmatter/nota_editorial}

% ===== Back Matter =====
\backmatter
\input{zh/frontmatter/glossario_historico}
\input{zh/frontmatter/nota_historica}
\input{zh/frontmatter/ficha_beta_diegetica}
\input{zh/frontmatter/backpage}

\end{document}
